\documentclass[../Main.tex]{subfiles}
\begin{document}
  \chapter{2023-2024学年微积分（一）（上）期末考试}

  \section{单项选择题(每小题 3 分，共 18 分)}
  \begin{enumerate}
    \item 下列函数在其定义域内无界的是【\quad】.

      \begin{tasks}[label=\textbf{\Alph*.}, label-width=1.5em, column-sep=1em](2)
        \task $x^{2}D(x)$（$D(x)$为Dirichlet函数）
        \task $\tan(\sin x)$
        \task $\frac{\sin x}{x}$
        \task 符号函数$\mathrm{sgn}(x)$
      \end{tasks}

    \item 曲线$y=\frac{1}{x}+\ln\left(1+\mathrm{e}^{x}\right)$的斜渐近线为【\quad】.

      \begin{tasks}[label=\textbf{\Alph*.}, label-width=1.5em, column-sep=1em](2)
        \task $y=-x$
        \task $y=-x+1$
        \task $y=x$
        \task $y=x+1$
      \end{tasks}

    \item 通解为$y=(C_{1}+C_{2} x)\mathrm{e}^{x}+x$（其中$C_{1},C_{2}$为任意常数）的微分方程为【\quad】.

      \begin{tasks}[label=\textbf{\Alph*.}, label-width=1.5em, column-sep=1em](2)
        \task $y''-y'=1$
        \task $y''-y=0$
        \task $y''-2y'+y=\mathrm{e}^{x}$
        \task $y''-2y'+y=x-2$
      \end{tasks}

    \item 设$f(x)$在$x=x_{0}$的某个邻域内有定义，则下列命题
      \begin{enumerate}
        \item 若$f'(x_{0})$存在，则$f(x)$在$x=x_{0}$处连续.

        \item 若$f'_{+}(x_{0}),f'_{-}(x_{0})$均存在，则$f(x)$在$x=x_{0}$处连续.

        \item 若$\lim_{x\to x_0^+}f'(x),\lim_{x\to x_0^-}f'(x)$均存在，则$f(x)$在$x
          =x_{0}$处连续.
      \end{enumerate}
      其中正确的个数是【\quad】.

      \begin{tasks}[label=\textbf{\Alph*.}, label-width=1.5em, column-sep=1em](2)
        \task $1$
        \task $2$
        \task $3$
        \task $0$
      \end{tasks}

    \item 设$f(x)$在$x=1$的某邻域内连续，且$\lim_{x\to0}\frac{\ln\left(f(x+1)+\mathrm{e}^{x^2}\right)}{x^{2}}
      =2$，则$x=1$是$f(x)$的【\quad】.

      \begin{tasks}[label=\textbf{\Alph*.}, label-width=1.5em, column-sep=1em](2)
        \task 驻点且为极大值点
        \task 驻点且为极小值点
        \task 不可导点
        \task 可导点但不是驻点
      \end{tasks}

    \item 设$M=\int_{-\frac{\pi}{2}}^{\frac{\pi}{2}}\cos^{4} x\sin^{3} x\,\mathrm{d}
      x$，$N=\int_{-\frac{\pi}{2}}^{\frac{\pi}{2}}\left(\cos^{3} x+\sin^{3} x\right
      )\,\mathrm{d}x$，$P=\int_{-\frac{\pi}{2}}^{\frac{\pi}{2}}\left(x\cos x-\sin^{2}
      x\right)\,\mathrm{d}x$，则【\quad】.

      \begin{tasks}[label=\textbf{\Alph*.}, label-width=1.5em, column-sep=1em](2)
        \task $M>N>P$
        \task $N>P>M$
        \task $N>M>P$
        \task $M>P>N$
      \end{tasks}
  \end{enumerate}

  \section{填空题(每小题 4 分，共 16 分)}
  \begin{enumerate}
    \setcounter{enumi}{6}
    \item 设曲线$y=f(x)$与$y=x^{2}-x+1$在点$(2,3)$处有公共切线，则$\lim_{n\to\infty}
      n\left[f\left(2-\frac{1}{n}\right)-3\right]=\underline{\hspace{2cm}}$.
      \vspace{0.5em}

    \item $y=f(x)$满足$y'=(1-y)y^{\alpha}(\alpha>0)$，若曲线$y=f(x)$的一个拐点为$\left
      (t,\frac{1}{2}\right)$，则$\alpha=\underline{\hspace{2cm}}$.
      \vspace{0.5em}

    \item 已知$f(x)$的一个原函数是$\mathrm{e}^{-x}$，则$\int_{-\frac{1}{2}}^{0}
      xf(2x)\,\mathrm{d}x=\underline{\hspace{3cm}}$.
      \vspace{0.5em}

    \item 位于曲线$y=x\mathrm{e}^{-x}(0\leq x<+\infty)$下方、$x$轴上方的无界图形的面积为$\underline
      {\hspace{3cm}}$.
      \vspace{0.5em}
  \end{enumerate}

  \section{计算题(每小题 7 分，共 42 分)}
  \begin{enumerate}
    \setcounter{enumi}{10}
    \item 设$y=y(x)$由方程$\mathrm{e}^{y}+6xy+x^{2}-1=0$确定，求$y'(0)$，$y''(0)$.
      \vspace{12em}

    \item 设$f(x)=\mathrm{e}^{x}\ln(1+x)-x\left(1+\frac{1}{2}x\right)$，当$x
      \to0$时，求$f(x)$的主部及阶数.
      \vspace{11em}

    \item 求极限$l=\lim_{x\to0}\frac{\mathrm{e}^{\sin x}}{x\sin x^{2}}\int_{0}
      ^{x}\sin(x-t)^{2} \mathrm{d}t$.
      \vspace{11em}

    \item 求极限$l=\lim_{n\to\infty}\frac{1}{n}\left[\ln\left(1+\frac{1}{n}\right
      )+\ln\left(1+\frac{2}{n}\right)+\cdots+\ln\left(1+\frac{n}{n}\right)\right]$.
      \vspace{11em}

    \item 设$f(x)=\int_{0}^{x}\frac{\sin t}{1-t}\,\mathrm{d}t$，求$I=\int_{0}^{1}
      (x-1)f(x)\,\mathrm{d}x$.
      \vspace{11em}

    \item 若曲线$y=f(x)$是$y''+2y'-3y=4\mathrm{e}^{x}$的一条积分曲线，此曲线过点$A
      (0,1)$，且在点$A(0,1)$处的切线的倾斜角为$\frac{3\pi}{4}$，求$f(x)$.
      \vspace{11em}
  \end{enumerate}

  \section{综合题(每小题 7 分，共 14 分)}
  \begin{enumerate}
    \setcounter{enumi}{16}
    \item 设$f(x)$在$(-\infty,+\infty)$上可导，其反函数存在为$g(x)$，若$\int
      _{0}^{f(x)}g(t)\,\mathrm{d}t+\int_{0}^{x}f(t)\,\mathrm{d}t=x\mathrm{e}^{x}-\mathrm{e}
      ^{x}+1$，求$f(x)$.
      \vspace{12em}

    \item 设函数$f(x)=ax+\frac{3}{2}bx^{2}$在区间$(0,1)$内恒大于$0$，其中$a,
      b$为未知常数.曲线$y=f(x)$与直线$x=1,y=0$所围成的区域$D$的面积为$2$. 求$a,b$的值，使得$D$绕$x$轴旋转一周得到的旋转体的体积最小，并求出此最小值.
      \vspace{13em}
  \end{enumerate}

  \section{证明题(每小题 5 分，共 10 分)}
  \begin{enumerate}
    \setcounter{enumi}{18}
    \item 设函数$f(x),g(x)$在区间$[0,a]$上连续且单调增加，其中$a>0$.证明
      \[
        a\int_{0}^{a}f(x)g(x)\,\mathrm{d}x\geq \int_{0}^{a}f(x)\,\mathrm{d}x\int_{0}^{a}
        g(x)\,\mathrm{d}x.
      \]
      \vspace{14em}

    \item 设$f(x)$在$[a,b]$上有二阶导数，且$f'(a)=f'(b)=0$. 证明：存在$\xi\in
      (a,b)$，使得
      \[
        (b-a)^{2}|f''(\xi)|\geq 4|f(b)-f(a)|.
      \]
      \vspace{14em}
  \end{enumerate}

  \chapter{2023-2024学年微积分（一）（上）期末考试参考答案}
  \section{单项选择题 (每小题 3 分，共 18 分)}
  \begin{enumerate}
    \item \textbf{Solution}. A.

      $x^{2}D(x)=
      \begin{cases}
        x^{2}, & x\in\mathbf{Q},                   \\
        0,     & x\in\mathbf{R}\setminus\mathbf{Q}
      \end{cases}$，显然$x^{2}D(x)$在其定义域内无界.

      $\forall\,x\in\mathbf{R}$，$|\sin x|\leq 1$，所以$|\tan(\sin x)|\leq \tan 1
      =\frac{\pi}{4}$，故$\tan(\sin x)$在其定义域内有界.

      $\lim_{x\to0}\frac{\sin x}{x}=1$，$\lim_{x\to\infty}\frac{\sin x}{x}=0$，所以$\frac{\sin
      x}{x}$在其定义域$(-\infty,0)\cup(0,+\infty)$内有界.

      $\mathrm{sgn}(x)=
      \begin{cases}
        -1, & x<0, \\
        0,  & x=0, \\
        1,  & x>0
      \end{cases}$，显然$\mathrm{sgn}(x)$在其定义域内有界.

    \item \textbf{Solution}. C.

      $\lim_{x\to+\infty}\frac{y}{x}=\lim_{x\to+\infty}\frac{\frac{1}{x}+\ln\left(1+\mathrm{e}^{x}\right)}{x}
      =\lim_{x\to+\infty}\frac{\ln\left(1+\mathrm{e}^{x}\right)}{x}=1$，

      $\lim_{x\to+\infty}(y-x)=\lim_{x\to+\infty}\left[\frac{1}{x}+\ln\left(1+\mathrm{e}
      ^{x}\right)-x\right]=\lim_{x\to+\infty}\ln\left(\frac{1+\mathrm{e}^{x}}{\mathrm{e}^{x}}
      \right)=\lim_{x\to+\infty}\ln\left(1+\mathrm{e}^{-x}\right)=0$，

      所以斜渐近线方程为$y=x$.

    \item \textbf{Solution}. D.

      由通解结构可以看出$(C_{1}+C_{2} x)\mathrm{e}^{x}$是一个齐次方程的通解，$y^{*}
      =x$是一个相应非齐次方程的特解.

      通解$(C_{1}+C_{2} x)\mathrm{e}^{x}$对应齐次方程的二重特征根$r=1$，所以齐次方程为$y
      ''-2y'+y=0$.

      将特解$y^{*}=x$代入非齐次方程验证可得D选项$y''-2y'+y=x-2$满足题意.

    \item \textbf{Solution}. B.

      若$f'(x_{0})$存在，即$\lim_{x\to x_0}\frac{f(x)-f(x_{0})}{x-x_{0}}$存在，所以$\lim
      _{x\to x_0}f(x)=f(x_{0})$，即$f(x)$在$x=x_{0}$处连续.

      同理，若$f'_{-}(x_{0})$存在，则$f(x)$在$x=x_{0}$处左连续；若$f'_{+}(x_{0})$存在，则$f
      (x)$在$x=x_{0}$处右连续，

      所以$f(x)$在$x=x_{0}$处连续.

      考虑函数$f(x)=
      \begin{cases}
        x^{2}, & x>0, \\
        0,     & x=0, \\
        -x+1,  & x<0
      \end{cases}$，则$\lim_{x\to0^+}f'(x)=\lim_{x\to0^+}2x=0$，$\lim_{x\to0^-}f'
      (x)=\lim_{x\to0^-}(-1)= -1$，

      但显然$f(x)$在$x=0$处不连续.

      所以有两个命题正确.

    \item \textbf{Solution}. B.
      \[
        \begin{aligned}
          \lim_{x\to0}\frac{\ln\left(f(x+1)+\mathrm{e}^{x^2}\right)}{x^2}= & \lim_{x\to0}\frac{f(x+1)+\mathrm{e}^{x^2}-1}{x^2}              \\
          =                                                                & \lim_{x\to0}\frac{f(x+1)}{x^2}+\lim_{x\to0}\frac{x^2}{x^2}= 2,
        \end{aligned}
      \]

      因此$\lim_{x\to0}\frac{f(x+1)}{x^{2}}=1$.由$f(x)$在$x=1$处的连续性易知$f(1)
      =0$.

      又由极限的保号性，当$x\to0$时，$\frac{f(x+1)}{x^{2}}>0$，所以$f(x+1)>0$. 即$x
      =1$是$f(x)$的极小值点.

      $\lim_{x\to 0}\frac{f(x+1)-f(1)}{x}=\lim_{x\to0}\frac{f(x+1)}{x^{2}}\cdot x
      =0$，所以$f'(1)=0$，即$x=1$是$f(x)$的驻点.

    \item \textbf{Solution}. C.

      $y=\cos^{4}\sin^{3} x$是奇函数，所以$M=\int_{-\frac{\pi}{2}}^{\frac{\pi}{2}}
      \cos^{4} x\sin^{3} x\,\mathrm{d}x=0$.

      $N=\int_{-\frac{\pi}{2}}^{\frac{\pi}{2}}\left(\cos^{3} x+\sin^{3} x\right)\,\mathrm{d}
      x=\int_{-\frac{\pi}{2}}^{\frac{\pi}{2}}\cos^{3} x\,\mathrm{d}x>0$.

      $P=\int_{-\frac{\pi}{2}}^{\frac{\pi}{2}}\left(x\cos x-\sin^{2} x\right)\,\mathrm{d}
      x=-\int_{-\frac{\pi}{2}}^{\frac{\pi}{2}}\sin^{2} x\,\mathrm{d}x<0$.

      所以$N>M>P$.
  \end{enumerate}

  \section{填空题(每小题 4 分，共 16 分)}
  \begin{enumerate}
    \setcounter{enumi}{6}
    \item \textbf{Solution}. $-3$.

      由题可知$f(2)=3$，$f'(2)=(2x-1)\bigg|_{x=2}=3$，

      所以$\lim_{n\to\infty}n\left[f\left(2-\frac{1}{n}\right)-3\right] = -\lim_{n\to\infty}
      \frac{f\left(2-\dfrac{1}{n}\right)-f(2)}{-\dfrac{1}{n}}= -f'(2) = -3$.

    \item \textbf{Solution}. $1$.

      方程$y'=(1-y)y^{\alpha}$两边对$x$求导得
      \[
        y''=-y'y^{\alpha}+(1-y)\alpha y^{\alpha-1}y'=\left[-y^{\alpha}+(1-y)\alpha
        y^{\alpha-1}\right]y'=\left[-y^{\alpha}+(1-y)\alpha y^{\alpha-1}\right](1
        -y)y^{\alpha}.
      \]

      由题意可知$y(t)=\frac{1}{2}$，$y''(t)=0$，代入得$\left[-\left(\frac{1}{2}\right
      )^{\alpha}+\frac{1}{2}\alpha \left(\frac{1}{2}\right)^{\alpha-1}\right]\left
      (\frac{1}{2}\right)^{\alpha+1}=0$，解得$\alpha=1$.

    \item \textbf{Solution}. $\frac{1}{4}$.

      由题可知，$f(x)=\left(\mathrm{e}^{-x}\right)'=-\mathrm{e}^{-x}$，所以
      \[
        \begin{aligned}
          \int_{-\frac{1}{2}}^{0}xf(2x)\,\mathrm{d}x = & \int_{-\frac{1}{2}}^{0}x\cdot\left(-\mathrm{e}^{-2x}\right)\,\mathrm{d}x                                                                                                             \\
          =                                          & \frac{1}{2}\int_{-\frac{1}{2}}^{0}x\,\mathrm{d}\mathrm{e}^{-2x}= \frac{1}{2}x\mathrm{e}^{-2x}\bigg|_{-\frac{1}{2}}^{0}-\frac{1}{2}\int_{-\frac{1}{2}}^{0}\mathrm{e}^{-2x}\,\mathrm{d}x \\
          =                                          & \frac{1}{4}\mathrm{e}+\frac{1}{4}\mathrm{e}^{-2x}\bigg|_{-\frac{1}{2}}^{0}= \frac{1}{4}.
        \end{aligned}
      \]

    \item \textbf{Solution}. $1$.

      由题可知，所求面积$A=$
      \[
        \int_{0}^{+\infty}x\mathrm{e}^{-x}\,\mathrm{d}x = -\int_{0}^{+\infty}x\,\mathrm{d}
        \mathrm{e}^{-x}= -x\mathrm{e}^{-x}\bigg|_{0}^{+\infty}+\int_{0}^{+\infty}
        \mathrm{e}^{-x}\,\mathrm{d}x = 1.
      \]
  \end{enumerate}

  \section{计算题(每小题 7 分，共 42 分)}
  \begin{enumerate}
    \setcounter{enumi}{10}
    \item \textbf{Solution}. 由方程得当$x=0$时，$y=0$.

      方程两边关于$x$求导得
      \[
        \mathrm{e}^{y}y'+6y+6xy'+2x=0.
      \]

      将$x=0$，$y=0$代入上式得$y'(0)=0$.

      方程两边关于$x$再次求导得
      \[
        \mathrm{e}^{y}(y')^{2}+\mathrm{e}^{y}y''+12y'+6xy''+2=0.
      \]

      将$x=0$，$y=0$，$y'(0)=0$代入上式得$y''(0)=-2$.

    \item \textbf{Solution}. 由Taylor公式可得
      \[
        \begin{aligned}
          f(x)=\mathrm{e}^{x}\ln(1+x)-x\left(1+\frac{1}{2}x\right) = & \left(1+x+\frac{1}{2}x^{2}+o(x^{2})\right)\left(x-\frac{1}{2}x^{2}+\frac{1}{3}x^{3}+o(x^{3})\right)-x-\frac{1}{2}x^{2} \\
          =                                                          & x-\frac{1}{2}x^{2}+\frac{1}{3}x^{3}+x^{2}-\frac{1}{2}x^{3}+\frac{1}{2}x^{3}+o(x^{3})-x-\frac{1}{2}x^{2}                \\
          =                                                          & \frac{1}{3}x^{3}+o(x^{3}).
        \end{aligned}
      \]

      所以$f(x)$的主部为$\frac{1}{3}x^{3}$，阶数为$3$.

    \item \textbf{Solution}. 令$u=x-t$，则$\int_{0}^{x}\sin(x-t)^{2} \mathrm{d}
      t = \int_{0}^{x}\sin u^{2}\,\mathrm{d}u$.

      故
      \[
        \begin{aligned}
          l = \lim_{x\to0}\frac{\mathrm{e}^{\sin x}}{x\sin x^2}\int_{0}^{x}\sin(x-t)^{2} \mathrm{d}t = & \lim_{x\to0}\frac{\displaystyle\int_{0}^{x}\sin u^2\,\mathrm{d}u}{x^3} \\
          =                                                                                            & \lim_{x\to0}\frac{\sin x^2}{3x^2}= \frac{1}{3}.
        \end{aligned}
      \]

    \item \textbf{Solution}. 由定积分的定义可得
      \[
        \begin{aligned}
          l= & \lim_{n\to\infty}\frac{1}{n}\sum_{k=1}^{n}\ln\left(1+\frac{k}{n}\right)                        \\
          =  & \int_{0}^{1}\ln(1+x)\,\mathrm{d}x = x\ln(1+x)\bigg|_{0}^{1}-\int_{0}^{1}\frac{x}{1+x}\,\mathrm{d}x \\
          =  & \ln 2 -\left(x-\ln(1+x)\right)\bigg|_{0}^{1}= \ln 2 - (1-\ln 2) = 2\ln 2 -1.
        \end{aligned}
      \]

    \item \textbf{Solution}. 由题可知$f(0)=0$，$f'(x)=\frac{\sin x}{1-x}$，所以
      \[
        \begin{aligned}
          I = \int_{0}^{1}(x-1) f(x)\,\mathrm{d}x = & \frac{1}{2}\int_{0}^{1}f(x)\,\mathrm{d}(x-1)^{2}                                                 \\
          =                                       & \left.\frac{1}{2}f(x)(x-1)^{2}\right|_{0}^{1}-\frac{1}{2}\int_{0}^{1}(x-1)^{2}f'(x)\,\mathrm{d}x \\
          =                                       & \frac{1}{2}\int_{0}^{1}(x-1)\sin x\,\mathrm{d}x = -\frac{1}{2}\int_{0}^{1}(x-1)\,\mathrm{d}\cos x  \\
          =                                       & -\frac{1}{2}\left[(x-1)\cos x\bigg|_{0}^{1}-\int_{0}^{1}\cos x\,\mathrm{d}x\right]               \\
          =                                       & \frac{\sin 1-1}{2}.
        \end{aligned}
      \]

    \item \textbf{Solution}. 齐次方程$y''+2y'-3y=0$的特征方程为$r^{2}+2r-3=0$，解得$r
      _{1}=1$，$r_{2}=-3$.

      对于$f(x)=4\mathrm{e}^{x}$，$\lambda = 1$是特征方程的单根，故可设特解为$y^{*}
      =Ax\mathrm{e}^{x}$，

      代入方程得$A(x+2)\mathrm{e}^{x}+2A(x+1)\mathrm{e}^{x}-3Ax\mathrm{e}^{x}=4\mathrm{e}
      ^{x}$，解得$A=1$，

      所以非齐次方程的通解为$y=C_{1}\mathrm{e}^{-3x}+C_{2}\mathrm{e}^{x}+x\mathrm{e}
      ^{x}$.

      由题可知$y(0)=1$，$y'(0)=\tan\frac{3\pi}{4}=-1$，代入得$\begin{cases}
        C_{1}+C_{2}=1,     \\
        -3C_{1}+C_{2}+1=-1
      \end{cases}$，解得$C_{1}=\frac{3}{4}$，$C_{2}=\frac{1}{4}$.

      因此$f(x)=\frac{3}{4}\mathrm{e}^{-3x}+\frac{1}{4}\mathrm{e}^{x}+x\mathrm{e}
      ^{x}$.
  \end{enumerate}

  \section{综合题(每小题 7 分，共 14 分)}
  \begin{enumerate}
    \setcounter{enumi}{16}
    \item \textbf{Solution}. 方程$\int_{0}^{f(x)}g(t)\,\mathrm{d}t+\int_{0}^{x}
      f(t)\,\mathrm{d}t=x\mathrm{e}^{x}-\mathrm{e}^{x}+1$两边关于$x$求导得$g(f(x))\cdot
      f'(x)+f(x)=x\mathrm{e}^{x}$.

      由反函数的性质，$g(f(x))=x$，所以上式即$xf'(x)+f(x)=x\mathrm{e}^{x}$.

      当$x=0$时，由题可知$f(0)=0$；当$x\neq0$时，上式可变形为
      \[
        f'(x)+\frac{1}{x}f(x)=\mathrm{e}^{x}.
      \]

      由一阶非齐次线性微分方程的通解公式得
      \[
        \begin{aligned}
          y = & \mathrm{e}^{-\int\frac{1}{x}\,\mathrm{d}x}\left(C+\int\mathrm{e}^{x}\mathrm{e}^{\int\frac{1}{x}\,\mathrm{d}x}\,\mathrm{d}x\right) \\
          =   & \frac{1}{x}\left(C+\int x\mathrm{e}^{x}\,\mathrm{d}x\right) = \frac{1}{x}\left(C+x\mathrm{e}^{x}-\mathrm{e}^{x}\right).
        \end{aligned}
      \]

      由$f(x)$在$x=0$处连续，所以$\lim_{x\to0}\frac{1}{x}\left(C+x\mathrm{e}^{x}-
      \mathrm{e}^{x}\right)=\lim_{x\to0}\left[\mathrm{e}^{x}+\frac{C-\mathrm{e}^{x}}{x}
      \right]=f(0)=0$，解得$C=1$.

      所以$f(x)=
      \begin{cases}
        0,                                                        & x=0,    \\
        \frac{1}{x}\left(1+x\mathrm{e}^{x}-\mathrm{e}^{x}\right), & x\neq0.
      \end{cases}$

    \item \textbf{Solution}. 由题可知$\int_{0}^{1}\left(ax+\frac{3}{2}bx^{2}
      \right)\,\mathrm{d}x=\frac{a}{2}+\frac{b}{2}=2$，所以$a=4-b$.

      旋转体的体积$V=$
      \[
        \begin{aligned}
          \int_{0}^{1}\pi f^{2}(x)\,\mathrm{d}x = & \pi\int_{0}^{1}\left[(4-b)x+\frac{3}{2}bx^{2}\right]^{2}\,\mathrm{d}x=\pi\int_{0}^{1}\left[\frac{9}{4}b^{2}x^{4}+3(4-b)bx^{3}+(4-b)^{2} x^{2}\right]\mathrm{d}x \\
          =                                     & \pi\left[\frac{9}{20}b^{2}+\frac{3}{4}(4-b)b+\frac{1}{3}(4-b)^{2}\right] = \left(\frac{1}{30}b^{2}+\frac{1}{3}b+\frac{16}{3}\right)\pi.
        \end{aligned}
      \]

      $V'=\left(\frac{1}{3}+\frac{1}{15}b\right)\pi$，令$V'=0$解得$b=-5$，此时$V''(-5)=\frac{\pi}{15}>0$，所以$V$的最小值存在，$V(-5)=\frac{9}{2}\pi$.

      故当$a=9$，$b=-5$时，$D$绕$x$轴旋转一周得到的旋转体的体积最小，此最小值为$\frac{9}{2}
      \pi$.
  \end{enumerate}

  \section{证明题(每小题 5 分，共 10 分)}
  \begin{enumerate}
    \setcounter{enumi}{18}
    \item \textbf{Proof}. 令$F(x)=x\int_{0}^{x}f(t)g(t)\,\mathrm{d}t-\int_{0}^{x}
      f(t)\,\mathrm{d}t\int_{0}^{x}g(t)\,\mathrm{d}t$，则$F(0)=0$.
      \[
        \begin{aligned}
          F'(x) = & \int_{0}^{x}f(t)g(t)\,\mathrm{d}t+x f(x)g(x)-f(x)\int_{0}^{x}g(t)\,\mathrm{d}t-g(x)\int_{0}^{x}f(t)\,\mathrm{d}t \\
          =       & \int_{0}^{x}\left[f(t)g(t)+f(x)g(x)-f(x)g(t)-g(x)f(t)\right]\mathrm{d}t                                    \\
          =       & \int_{0}^{x}\left[f(x)-f(t)\right]\cdot\left[g(x)-g(t)\right]\mathrm{d}t.
        \end{aligned}
      \]

      由题可知$f(x),g(x)$在$[0,a]$上单调增加，所以$\forall\,t\in[0,x]$，$f(x)-f(t
      )\geq0$，$g(x)-g(t)\geq0$，

      故$\left[f(x)-f(t)\right]\cdot\left[g(x)-g(t)\right]\geq0$，因此$F'(x)\geq0$，$F
      (x)$在$[0,a]$上单调增加.

      所以$F(x)\geq F(0)=0$，即
      \[
        a\int_{0}^{a}f(x)g(x)\,\mathrm{d}x\geq \int_{0}^{a}f(x)\,\mathrm{d}x\int_{0}^{a}
        g(x)\,\mathrm{d}x.
      \]

    \item \textbf{Proof}. 由Taylor公式，将$f\left(\frac{a+b}{2}\right)$分别在$x
      =a$和$x=b$处展开得
      \[
        \begin{gathered}
          f\left(\frac{a+b}{2}\right)=f(a)+f'(a)\left(\frac{a+b}{2}-a\right)+\frac{1}{2}f''(\xi_1)\left(\frac{a+b}{2}-a\right)^{2},\\
          f\left(\frac{a+b}{2}\right)=f(b)+f'(b)\left(\frac{a+b}{2}-b\right)+\frac{1}{2}f''(\xi_2)\left(\frac{a+b}{2}-b\right)^{2}.
        \end{gathered}
      \]

      因$f'(a)=f'(b)=0$，所以上式即
      \[
        \begin{gathered}
          f\left(\frac{a+b}{2}\right)=f(a)+\frac{1}{8}f''(\xi_1)(b-a)^{2},\\
          f\left(\frac{a+b}{2}\right)=f(b)+\frac{1}{8}f''(\xi_2)(b-a)^{2},
        \end{gathered}
      \]

      其中$\xi_{1}\in\left(a,\frac{a+b}{2}\right)$，$\xi_{2}\in\left(\frac{a+b}{2}
      ,b\right)$.

      两式相减得
      \[
        0=f\left(\frac{a+b}{2}\right)-f\left(\frac{a+b}{2}\right)=f(a)-f(b)+\frac{1}{8}
        \left[f''(\xi_{1})-f''(\xi_{2})\right](b-a)^{2},
      \]

      所以
      \[
        |f(b)-f(a)|=\frac{(b-a)^{2}}{8}\left|f''(\xi_{1})-f''(\xi_{2})\right|\leq
        \frac{(b-a)^{2}}{8}\left(|f''(\xi_{1})|+|f''(\xi_{2})|\right).
      \]

      取$|f''(\xi)|=\max\{|f''(\xi_{1})|,|f''(\xi_{2})|\}$，

      则$|f(b)-f(a)|\leq \frac{(b-a)^{2}}{8}\cdot 2|f''(\xi)|$，即$(b-a)^{2}|f''(
      \xi)|\geq 4|f(b)-f(a)|$.
  \end{enumerate}
\end{document}





